\section{Model adaptations and training}
\label{app:protocol}
The library varies both the representation of a field and how coarse data
enter learning. We first identify the retained published components and our
adaptations, then give the training settings. The additional baseline
descriptions use the same distinction between a source idea and a complete
published method.

\subsection{Model sources, adaptations, and rationale}
\label{app:adaptations}
Each library member combines a field representation with a training and inference
procedure. We describe M1 to M9 individually below, identifying the published
components retained, the changes made for parameter to field prediction, and
the reason for including each model. Several members combine ideas from more
than one paper. The resulting library is an integration of adapted models,
not nine independent reproductions or nine newly proposed architectures.
All accept parameters alone at deployment. Appendix~\ref{app:experts} gives
their training budgets and data use. The rationales below motivate the library's
composition, while the experiments assess its predictions. They do not isolate
the effect of every architectural change.

\paragraph{M1: FiLM FNO transfer.}
The starting components are FNO's learned Fourier mixing \citep{li2021fno},
FiLM's featurewise affine conditioning \citep{perez2018}, and the LF pretraining
followed by HF finetuning strategy of \citet{lyu2023}. We retain spectral
convolutions with local residual paths and adapt the input interface to a
simulation parameter vector. The spatial input contains coordinates only.
At each block, a parameter network generates a scale and shift applied after
group normalization, so the parameters influence features throughout the
network. Related parameter conditioning for PDE emulation is studied by
\citet{shokar2025}. Our model differs from our parameter concatenation FNO references, which
supply parameters as constant spatial channels at the input. The same weights
are first trained on LF fields and then finetuned on HF fields. This combination
tests whether repeated parameter conditioning and a coarse initialization
provide an effective way to generate globally coupled fields.

\paragraph{M2: All pairs FNO.}
M2 retains M1's FNO and FiLM components \citep{li2021fno,perez2018} and adapts
the idea of pair enumeration from Poseidon \citep{herde2024}. Poseidon uses
pairs of states for temporal operator learning. Here the pairs index source
and target fidelities. We append their normalized labels to the parameter
vector, pool the corresponding target fields on a common working grid, and
train one conditional FNO before HF finetuning. The source label specifies a
training pair, but no source field enters the network. Poseidon's architecture,
pretrained weights, and temporal operator are not used. For ERA5, exact input
joins retain each target field with its own parameters. The purpose is to let
one model share information across fidelity tasks before specializing to the
fine field. Pair enumeration also repeats targets and changes optimization
cost, so its effect cannot be attributed to fidelity conditioning alone.

\paragraph{M3: ConvNeXt transfer.}
ConvNeXt supplies the depthwise spatial convolution, channel expansion, and
residual block structure \citep{liu2022convnext}. We place these blocks in an
U shaped encoder and decoder with skip connections \citep{ronneberger2015} and
a dense field output, adapting the image classification backbone to field
generation. The model starts from
coordinates and replaces the usual block normalization with group normalization
and FiLM conditioning on the simulation parameters \citep{perez2018}. Three
downsampling stages and a bottleneck process the field at several resolutions,
with corresponding decoder stages restoring the target grid. LF pretraining
and HF finetuning follow the transfer principle of \citet{lyu2023}. The
rationale is to provide local spatial processing alongside M1's Fourier mixing,
while retaining the same type of parameter conditioning and transfer schedule.

\paragraph{M4: Distribution FNO.}
FIRE conditions a fine correction on summaries of a coarse predictive
distribution using tabular foundation models and in context inference
\citep{yu2026fire}. M4 retains the distribution conditioning idea and replaces
those models with FNOs trained on spatial fields \citep{li2021fno}. Five LF
members predict a coarse field at the query parameters. Their pointwise mean,
standard deviation, and 10th, 50th, and 90th percentiles become spatial input
channels for a residual FNO, together with coordinates and broadcast parameters.
The residual is added to the coarse mean to predict the fine field. Both stages
are trained by gradient optimization, so this implementation does not retain
FIRE's inference without model retraining. The ensemble summaries are features
of prediction variability, not a guarantee of accurate uncertainty estimates. The
rationale is to expose coarse disagreement as well as the mean to the
correction model. This recipe has no out of fold construction and is distinct
from the coarse ensembles used by M8 and M9.

\paragraph{M5: Wavelet transfer.}
WNO motivates learning in wavelet coordinates \citep{tripura2022wno}. We
replace M1's Fourier mixing with a custom single level db4 wavelet transform,
while retaining its coordinate input, FiLM conditioning \citep{perez2018},
local residual paths, and LF to HF transfer schedule \citep{lyu2023}. Within
each wavelet subband, learned weights mix channels at coefficient positions
inside a capped leading block. Other positions pass through unchanged before
the inverse transform. These indices describe positions within a subband,
not a Fourier frequency cutoff. The implementation also handles rectangular
grids, odd grid sizes, and one dimensional grids. It is a custom wavelet
representation rather than an exact reproduction of the published WNO.
Its inclusion tests whether spatially localized wavelet features provide useful
predictions alongside the spectral and convolutional transfer models.

\paragraph{M6: Retrieved field DeepONet.}
DeepONet supplies the branch and trunk representation \citep{lu2021}, while
the multifidelity use of a coarse prediction as a reference is motivated by
\citet{lu2022mf}. We use a Cartesian product DeepONet whose branch receives
the parameters concatenated with a stored LF training field and whose trunk
receives target grid coordinates. A nearest neighbour search in standardized
parameter space selects that field. The network learns a residual relative to
its interpolated version, which is then added back to form the fine prediction.
Retrieval is the deployment adaptation: a new query requires parameters only
and no new coarse solve. The rationale is to exploit a nearby physical field
without first learning a separate coarse surrogate. If an exact LF match exists
at a training input, retrieval supplies that match, whereas an unseen query
uses a neighboring input. This difference in reference quality remains a
limitation. The implemented residual model is also distinct from the composite
multifidelity architecture of \citet{howard2023}.

\paragraph{M7: POD GP.}
Basis emulation represents a large field through a small set of coefficients
\citep{higdon2008}. M7 retains this principle, constructing a POD basis by
singular value decomposition of centered HF training fields. A Gaussian process
maps parameters to each retained coefficient, and the predicted coefficients
are projected back through the basis and added to the training mean field.
Our implementation uses an RBF kernel with automatic relevance determination
and shares covariance hyperparameters across coefficient outputs. It does not
implement the full Bayesian hierarchy of the cited model. M7 is deliberately
fine only: no coarse field enters basis construction, fitting, or inference.
It supplies a compact alternative when the field varies mainly in a few
directions or when learning a coarse correction is less effective than learning
the fine coefficients directly.

\paragraph{M8: Slice attention corrector.}
Transolver motivates grouping spatial cells into learned slice tokens
\citep{wu2024transolver}. We retain pooled spatial summaries but change the
attention operation. Original Transolver normalizes each cell's assignments
across slices, applies attention between slice tokens, and maps back using
the assignments. Our pooling weights instead normalize across cells, and
cell queries attend directly to the pooled slice tokens. The model receives
the current field, coordinates with sinusoidal features, and FiLM conditioning
on the parameters \citep{perez2018}. A coarse surrogate ensemble supplies its
initial state, with out of fold predictions for training cases. A head
initialized to zero predicts additive corrections through six damped updates
with shared weights. Training penalizes spatial squared error and Fourier
magnitude discrepancy at each update, plus a nonzero update when started at
the target field. This encourages an equilibrium without guaranteeing
contraction. The adaptation tests whether pooled global interactions can
refine a predicted coarse field using parameters alone at deployment. Its
attention and iterative recipe differ from the published Transolver.

\paragraph{M9: ConvNeXt corrector.}
M9 retains ConvNeXt's spatial and channel mixing blocks \citep{liu2022convnext}
and adapts them to the same correction task as M8. It uses a custom U shaped
encoder and decoder with skip connections \citep{ronneberger2015}, three
downsampling stages and a bottleneck,
and FiLM parameter conditioning \citep{perez2018}. Unlike M3, which generates
a field from coordinates and parameters after transfer training, M9 receives
the current coarse estimate as a spatial input and outputs an additive update.
Its head is initialized to zero, making the initial refinement the identity.
The coarse ensemble construction, six damped updates, shared weights, and
spatial, spectral, and target state penalties follow the M8 recipe. The rationale
is to offer local processing at several resolutions as an alternative to pooled
attention when correcting the same kind of coarse prediction. This is a
ConvNeXt based iterative corrector, not the original classification model.

The additional baseline adaptations below use the same distinction between a
source idea and a complete published method. In particular, B1, B5, and B8 omit
central probabilistic components of co kriging, DMFAL, and IFC. Their scores
describe the implemented references, not faithful reproductions of those
complete methods. Common evaluation cases make the scores comparable while
leaving architecture and resource effects coupled.

\paragraph{B1, B2, and B11: classical multifidelity field references.}
B1 applies an autoregressive discrepancy idea \citep{kennedy2000} to POD GP field
models. It fits one scalar coupling coefficient in field space, constructs a
separate POD GP for residuals, and combines posterior means at deployment.
This is a plug in approximation rather than joint probabilistic co kriging.
B2 introduces nonlinear dependence on LF coefficients \citep{perdikaris2017}.
Its HF coefficient kernel combines input dependent coupling, an LF coefficient
kernel, and a discrepancy kernel. It shares covariance hyperparameters across
HF outputs and propagates marginal LF uncertainty through Monte Carlo samples.
B11 is a cheaper affine alternative inspired by multifidelity surrogate work
\citep{zhang2018}. One slope and one intercept are fitted jointly over HF training cases and
pixels, using observed LF fields when paired and emulated LF fields otherwise.
They are then applied to the LF POD GP prediction.
It does not fit an affine mapping separately for each mode. Paired training rows
can use observed LF fields, whereas deployment uses predicted LF fields.

\paragraph{B3, B4, and B5: neural fusion references.}
B3 uses the composite correlation idea of \citet{meng2020}. A tanh MLP maps
parameters and coordinates to an LF value. This stage is frozen before linear
and nonlinear HF networks are trained on the original inputs and LF prediction,
with weight decay on the nonlinear branch. B4 follows the composite DeepONet
idea \citep{howard2023}, predicting fixed coarse sensors with an LF DeepONet and
feeding parameters plus sensors to linear and nonlinear HF DeepONets. The LF
stage is frozen and the HF branches are trained jointly in a subsequent stage. B5 takes hierarchical
latent transfer as inspiration from DMFAL \citep{li2022dmfal}. It uses two encoders
and linear field decoders, with the second encoder receiving parameters and the
first latent vector. Joint LF and HF squared losses replace the probabilistic
objective. Both the variational last layer and active acquisition are omitted.
The name latent MF network reflects these substantial departures.

\paragraph{B6 and B7: alternative decoder and probabilistic models.}
B6 uses NOMAD's nonlinear decoding principle \citep{seidman2022}. A parameter MLP
produces a latent vector and a coordinate conditioned MLP decodes the field.
The model is implemented independently, with multifidelity learning supplied by
LF pretraining and HF finetuning. B7 directly imports MFRNP's upstream
\texttt{MultiFidelityModel} \citep{niu2024}. Its adapter changes data loading,
normalization, resizing, and training. A scalar mean and standard deviation from
the first fidelity normalize inputs, while outputs are scaled per fidelity.
The longest grid side is capped at 256, rectangular grids are supported, and
bilinear antialiased resizing replaces the upstream transform. Batch size is
limited by the smallest fidelity, with independent shuffling across fidelities.
When HF training has at most ten examples, all rows are retained and the context
fraction range changes from $[0.20,0.25]$ to $[0.40,0.60]$. Hidden width is 128,
latent width is 64, and there are three hidden layers. Adam uses learning rate
and epsilon of $10^{-3}$. HF likelihood receives weight two and LF likelihood
weight one. Latent summaries are saved from the final training batch. These are
protocol adaptations of an upstream model, rather than a reproduction of its
published experiment.

\paragraph{B8, B9, and B10: fidelity and transfer FNO references.}
B8 forms latent spatial channels with an FNO and mixes them with weights predicted
from $[\ell,\ell^2]$, the fidelity label and its square. These two features
condition the spatial basis weights on fidelity. It uses maximum absolute scaling
per fidelity, the nearest registered scaler for an intermediate fidelity, LF
warm up, and joint LF and HF finetuning. IFC motivates fidelity dependent
representation \citep{li2022ifc}, but its neural ODE and Gaussian process are
absent. The descriptive name fidelity basis FNO makes this distinction explicit.
B9 and B10 are two archived entries for parameter concatenation FNO transfer
\citep{lyu2023}. In the inspected release, their model files are byte identical
and their training scripts differ only in repository paths and reported
identifiers. The labels FNO transfer I and II distinguish saved result entries,
not independent baseline families. Their different saved scores remain
unexplained by the inspected source. Both are retained to make the archived
comparison auditable, but neither their number nor their score difference is
evidence for architectural diversity or a reproducible recipe difference.
For ERA5 and the four additional reaction diffusion datasets, the identical
recipe is fitted once and its predictions are explicitly reused for B10,
with no additional training cost.


\subsection{Training settings and data use}
\label{app:experts}
Table~\ref{tab:training} summarizes the library schedules. For M1, M2, M3,
and M5, AdamW uses batch size at most 16, weight decay $10^{-5}$, gradient
clipping at 1, and cosine decay to $10^{-6}$. Initial learning rates are
$10^{-3}$ for pretraining and $3\times10^{-4}$ for HF finetuning. The FNO
recipe uses width 64, four blocks, and up to 12 Fourier modes per direction,
subject to the grid cap. M4 also uses AdamW but starts both its LF members
and residual stage at $10^{-3}$. M6 uses width 256 and five affine layers of
that width in each branch and trunk specification, including the latent output,
with ReLU and Adam at $10^{-3}$. Detailed grid dependent definitions and fitted
hyperparameters are retained with the model implementations and metadata.

\begin{table}[h!]
\centering\small\setlength{\tabcolsep}{4pt}
\caption{\textbf{Library training schedules.} Epochs count passes through the
relevant training table. Optimizer steps and repeated refinement steps are
different quantities. Each final surrogate uses a single seed, with internal
coarse ensemble initializations prescribed by its recipe.}
\label{tab:training}
\begin{tabular}{p{.16\linewidth}p{.72\linewidth}}
\toprule
Models & Nominal training procedure \\
\midrule
M1, M3, M5 & 2,500 LF epochs, followed by 2,500 HF epochs. \\
M2 & 2,500 epochs over pooled fidelity pairs, followed by 1,250 HF epochs. \\
M4 & Five LF FNOs and one HF residual FNO, each trained for 2,500 epochs. \\
M6 & 125,000 Adam iterations on HF residual targets. \\
M7 & HF basis construction and coefficient GP fitting, without an epoch schedule. \\
M8, M9 & Coarse members request 30,000 optimizer steps, rounded to whole epochs.
Each corrector uses 6,000 optimizer steps and six refinement updates per prediction. \\
\bottomrule
\end{tabular}
\end{table}

M1, M3, M4, M5, and M6 use the lowest registered LF level. M2 pools pairs
across registered fidelities, M8 and M9 predict the finest registered LF
level, and M7 uses HF fields only. ERA5 consequently uses level one for the
transfer and retrieval recipes and level eight for the iterative correctors.
For M8 and M9, five LF folds each contain four members: plain and
heteroscedastic variants with two initializations each. Training cases receive
the mean from their held out fold. All query cases use the four members of
fixed fold zero. Corrector checkpoints are selected using training only
validation. M7's basis uses only HF base training fields. Its original estimator
state was not saved, so the common prediction export reconstructs that training
only fit and restores target scaling.

Baselines use the settings of the evaluated benchmark implementations, with
the adaptations above. Training, validation allocations, and computation vary
across recipes. For example, on Darcy, B3 to B5 use 360 fine rows for fitting
and 40 for internal validation, while B1, B2, and the transfer recipe fit
400 rows. Both allocations consume 400 fine labels before ensemble fitting.
The total fine data budget is $N_H^{\rm base}+K$, plus any additional HF labels
used for tuning, where $N_H^{\rm base}$ counts base training labels and $K$
counts ensemble fitting examples. Final evaluation labels are separate.
Computation includes every candidate and coarse member trained or evaluated,
even if its final weight is zero. Some timing records measure prediction
export rather than training, preventing a common training cost estimate.
Thus common evaluation cases do not establish equal resources or isolate a
multifidelity advantage over a matched fine only ensemble.

